\documentclass[11pt,a4paper]{article}
\usepackage[utf8]{inputenc}
\usepackage[T1]{fontenc}
\usepackage[spanish,es-nodecimaldot]{babel}
\usepackage{amsmath,amssymb,bm,booktabs,array,graphicx,geometry}
\usepackage[colorlinks=true,allcolors=blue]{hyperref}
\geometry{margin=24mm}
\title{Respuesta angular calibrada y PEB del receptor reorientable\\\large Lectura crítica de Bastiaens et al. y aplicación al PDA100A2}
\author{Proyecto Cambridge}
\date{27 de septiembre de 2026}
\begin{document}
\maketitle
\begin{abstract}
Se estudia el artículo de Bastiaens et al., IEEE Access 2020, sobre la respuesta angular de los Thorlabs PDA100A2 y PDA36A2. Se separan sus ajustes experimentales de las hipótesis adicionales del sistema Cambridge: un LED fijo, un PD reorientable, ruido óptico equivalente y ganancia calibrada. La FIM se deriva para una respuesta axisimétrica general $h(s)$, no exclusivamente para dos dispositivos. Los modelos SQ, SQapprox y Exp se implementan como casos documentados, conservando el coseno-potencia como familia independiente. Se examinan los límites angulares, la identificabilidad, la incertidumbre de pose y la relevancia de esta generalización para un artículo de revista.
\end{abstract}
\tableofcontents
\clearpage

\section{Qué investigó realmente el artículo}
\emph{Impact of a Photodiode's Angular Characteristics on RSS-Based VLP Accuracy} \cite{bastiaens2020} caracteriza dos receptores comerciales y muestra que una respuesta angular mejor descrita puede reducir el desajuste de modelos RSS. El experimento de localización utiliza cuatro LEDs COB y un receptor que se desplaza en un plano aproximadamente 3 m por debajo; no es el sistema de un LED fijo y un PD reorientable de Cambridge.

Se distinguen dos conjuntos de datos: \emph{Point}, medido inclinando el PD bajo un único LED a distancia 1.25 m, y \emph{Plane}, obtenido de medidas espacialmente distribuidas bajo los cuatro LEDs. En Point se compensan traslaciones del mecanismo para mantener el centro óptico en la misma posición. Los ajustes angulares se apoyan en Plane; los autores lo consideran más fiable que la operación manual del mecanismo de Point.

El artículo compara modelos de aceptación normalizada, ajustes empíricos RSS--distancia, model-based fingerprinting (MBF) y trilateración. Para PDA100A2, la Tabla 4 da un percentil 90 de error MBF de 54.58 cm con coseno simple, 10.06 cm con coseno de orden 1.9 y 9.22 cm con SQ. La ganancia SQ frente al coseno ajustado es 0.84 cm (8.4\%), no la ganancia del 83.1\% frente al coseno simple. Estos números pertenecen al montaje y algoritmo de ese artículo; no se transfieren a nuestro error 3D ni a nuestro PEB.

\subsection{Lo que no elimina la propuesta}
La afirmación de evitar ajustes arbitrarios RSS--distancia \emph{por LED} no significa que el sistema funcione sin calibración. El artículo compensa tilt de LEDs, calibra sus potencias, utiliza un patrón de irradiación tabulado para MBF y considera un factor de desajuste espectral. Sus propias secciones III-A1 explicitan esas calibraciones. Las superficies negras reducen multitrayectoria; no se demuestra que un modelo angular resuelva reflejos, movimiento o deriva radiométrica.

\section{Dispositivos, parámetros y trazabilidad}
Los valores siguientes proceden de la Tabla 1 y de la sección III-A2. Son parámetros publicados redondeados, no una recalibración de nuestra unidad física.
\begin{center}
\begin{tabular}{lrr}
\toprule Parámetro & PDA100A2 & PDA36A2 \\
\midrule
Área activa nominal [mm$^2$] & 75.4 & 13 \\
Ganancia electrónica usada en localización [V/A] & $4.75\times10^4$ & $4.75\times10^5$ \\
Orden coseno ajustado $m_R$ & 1.9 & 0.98 \\
Ángulo SQ $\psi_{3\mathrm{dB}}$ [rad] & 0.61 & 0.79 \\
Ángulo Exp $\psi_{1/2}$ [rad] & 0.80 & 1.06 \\
Exponente Exp $S_R$ & 2.39 & 2.30 \\
\bottomrule
\end{tabular}
\end{center}
El PDA100A2 es el preset principal: coincide con el área de $75.4\ \mathrm{mm}^2$ del proyecto y con el orden 1.9 utilizado previamente. El PDA36A2 sirve como segundo caso de respuesta. En comparaciones de \emph{forma angular}, todos los modelos conservan la misma área, potencia y ruido óptico de Cambridge; el área de referencia del fabricante se guarda aparte. No se presenta esa comparación normalizada como una comparación completa de prestaciones de dos dispositivos.

La ganancia transimpedancia se guarda como metadato bibliográfico y no se multiplica en una media ya expresada en vatios ópticos equivalentes. El modelo normalizado incluye la respuesta angular efectiva del receptor; no debe multiplicarse otra vez por $\cos\psi$. En el dominio eléctrico habría que convertir consistentemente tanto media como covarianza usando ganancia y responsividad normal.

\subsection{Tres niveles distintos de afirmación}
\begin{enumerate}
\item \textbf{Publicado}: forma de las ecuaciones, coeficientes redondeados y resultados del artículo.
\item \textbf{Implementado}: evaluar esas funciones y derivarlas analíticamente bajo hipótesis explícitas de soporte angular y ruido.
\item \textbf{Validado en nuestro hardware}: requeriría medidas independientes del PD concreto, espectro, ganancia, temperatura, paquete y ruido. Este tercer nivel no se infiere automáticamente de los dos primeros.
\end{enumerate}
No se han inventado puntos de las curvas medidas. Tampoco se recalculó $R^2$ sin los datos originales. Hay una discrepancia menor en el artículo: el texto atribuye a SQapprox $R^2=0.983/0.982$, mientras su Tabla 2 muestra $0.989/0.983$ para PDA36A2/PDA100A2. No se usa esa inconsistencia para ajustar coeficientes. Los $R^2$ de SQ de esa tabla son 0.991/0.993.

\section{Funciones de respuesta y significado de FOV}
Sea $s=\cos\psi$ con incidencia física $\psi\in[0,\pi]$, y sea $h(s)$ la aceptación normalizada. En el hemisferio frontal se aplica además una apertura externa explícita $\Psi_{\rm ext}\leq90^\circ$.

\subsection{Familia general coseno-potencia}
\begin{equation}
 h(s)=s^{m_R},\qquad h'(s)=m_Rs^{m_R-1},\quad
 s>0,\ \psi\leq\Psi_{\rm ext}.
\end{equation}
Es la familia histórica. $m_R=1$ sigue siendo el valor general por defecto. Un orden 5 sería un control hipotético, no el ajuste que el artículo da para PDA100A2. La comparación principal de K=5 es con $m_R=1.9$.

\subsection{SQ, ecuación (5) del artículo}
\begin{equation}
 a=\frac{1-1/\sqrt2}{\psi_{3\mathrm{dB}}^2},\qquad
 h(s)=1-a(\arccos s)^2\quad\text{dentro del soporte},
\end{equation}
\begin{equation}
 \psi_C=\min\{\Psi_{\rm ext},\pi/2,1/\sqrt a\},\qquad h=0\quad\text{fuera}.
\end{equation}
Para PDA100A2 el cero natural es $64.580^\circ$; para PDA36A2 es $83.636^\circ$. \textbf{SQ no elimina todo corte}: llega a cero en un ángulo finito y luego permanece en cero. Si se usa el cero natural es continuo, pero su primera derivada cambia allí. Si una apertura externa corta antes, puede existir un salto positivo a cero.

El parámetro $\psi_{3\mathrm{dB}}$ corresponde a $h=1/\sqrt2$, debido a una caída de 3 dB en \emph{potencia eléctrica}, proporcional al cuadrado de amplitud. No es el ángulo de media respuesta $h=1/2$. Los parámetros del artículo están en radianes; sustituir 0.61 por grados cambia completamente el modelo.

La derivada requerida por la FIM es:
\begin{equation}
 h'(s)=2a\frac{\arccos s}{\sqrt{1-s^2}}=2a\frac\psi{\sin\psi},\qquad
 h'(1)=2a.
\label{eq:sqderivative}
\end{equation}
La singularidad aparente en incidencia normal es removible. El código usa su límite/serie, no divide cero entre cero.

\subsection{SQapprox, ecuación (6)}
\begin{equation}
 b=\frac{2-\sqrt2}{\psi_{3\mathrm{dB}}^2}=2a,\qquad
 h(s)=1-b+bs,\qquad h'(s)=b.
\end{equation}
Si $b>1$, el cero frontal se encuentra en $\psi_C'=\arccos(1-1/b)$. Para PDA100A2 vale $68.606^\circ$. En PDA36A2 el cero de la función afín queda fuera del hemisferio frontal; se conserva el límite físico externo de $90^\circ$. Por tanto SQ y SQapprox no poseen necesariamente el mismo soporte.

El nombre ``approx'' es relevante: el artículo obtiene esta función aproximando $\psi^2$ por $2(1-\cos\psi)$. No debe etiquetarse como la misma curva SQ. Es algebraicamente útil, pero su ajuste puede ser peor.

\subsection{Exp, ecuación (3)}
\begin{equation}
 h(s)=\exp[-\lambda(\arccos s)^{S_R}],\qquad
 \lambda=\frac{\ln2}{\psi_{1/2}^{S_R}},
\end{equation}
\begin{equation}
 h'(s)=\lambda S_R\psi^{S_R-2}\frac\psi{\sin\psi}\,h(s).
\end{equation}
Los exponentes publicados son mayores que dos, por lo que $h'(1)=0$. Exp no llega a cero a un ángulo finito; el soporte frontal externo sigue siendo una hipótesis adicional. El artículo observa divergencia creciente de Exp respecto a Point para $\psi>1.3$ rad. No se toma su cola como medida fiable de respuesta real. Los experimentos usan un límite externo de $90^\circ$ explícito y lo identifican como hipótesis de hemisferio frontal, no como una FOV certificada del fabricante.

\subsection{Modelos no reproducidos como presets}
El artículo considera $\cos(a\psi)$ y correcciones Fresnel. No se ha inventado el coeficiente $a$ del ajuste MUSIC ni una tabla angular experimental. El valor 4.22 mencionado procede del índice óptico del silicio ponderado espectralmente; no es una reflectancia admisible superior a uno. Una extensión Fresnel exige explicitar índice complejo/espectro, polarización y normalización respecto a incidencia normal. Los presets principales reproducen las ecuaciones SQ, SQapprox y Exp con parámetros publicados, sin añadir una corrección no verificada ni contar pérdidas dos veces.

\section{Canal general del sistema Cambridge}
Se define:
\begin{equation}
 \mathbf v=\mathbf t-\mathbf r,\quad d=\|\mathbf v\|,\quad \mathbf u=\mathbf v/d,
 \quad s_i=\mathbf n_i^{\mathsf T}\mathbf u,\quad \beta=\frac{Cf_T(\mathbf u)}{d^2}.
\end{equation}
$\mathbf u$ apunta del PD al LED. En una región regular de soporte fijo:
\begin{equation}
 \mu_i(\mathbf r)=\beta(\mathbf r)h(s_i),\qquad
 \bar Y_i\sim\mathcal N(\mu_i,\sigma^2/N_i).
\label{eq:generalchannel}
\end{equation}
$h$ ya es la respuesta angular efectiva total. El LED no se reorienta y el centro óptico no se desplaza. $f_T$ puede ser una respuesta emisora positiva calibrada, no necesariamente Lambertiana. El caso Lambertiano usa $f_T=c^{m_T}$, $c=(-\mathbf n_t)^{\mathsf T}\mathbf u$, y $C=P_t(m_T+1)A_{\rm PD}T_sg/(2\pi)$.

\section{Derivación del Jacobiano y la FIM}
Con $P_u=I-\mathbf u\mathbf u^{\mathsf T}$:
\begin{equation}
 \nabla d=-\mathbf u,\qquad D_{\mathbf r}\mathbf u=-P_u/d,\qquad
 \nabla s_i=-\frac{\mathbf n_i-s_i\mathbf u}{d}.
\end{equation}
El gradiente de amplitud es:
\begin{equation}
 \nabla\beta=\frac{C}{d^3}\left[2f_T\mathbf u-P_u\nabla_{\mathbf u}f_T\right].
\end{equation}
Por la regla del producto:
\begin{equation}
 \boxed{\nabla\mu_i=h(s_i)\nabla\beta-
 \frac\beta d h'(s_i)(\mathbf n_i-s_i\mathbf u).}
\label{eq:gradient}
\end{equation}
Para emisión Lambertiana se obtiene la expresión equivalente:
\begin{equation}
 \nabla\mu_i=\frac\beta d\left[-h'_i\mathbf n_i+
 ((m_T+2)h_i+s_i h'_i)\mathbf u-\frac{m_T h_i}{c}(-\mathbf n_t)\right].
\end{equation}
Al sustituir $h=s^{m_R}$ se recupera el gradiente generalizado anterior con coeficiente radial $(m_T+m_R+2)s^{m_R}$. Para SQ se usa \eqref{eq:sqderivative}; no basta con cambiar la media y conservar el gradiente del coseno.

Con $J$ de filas $\nabla\mu_i^{\mathsf T}$ y $W=\operatorname{diag}(N_i/\sigma^2)$, la puntuación es $J^{\mathsf T}W(\bar{\mathbf Y}-\bm\mu)$. Su segundo momento da \cite{kay}:
\begin{equation}
 I_r=J^{\mathsf T}WJ,\qquad
 \mathrm{PEB}=\sqrt{\operatorname{tr}(I_r^{-1})}.
\end{equation}
Se utiliza SVD de $W^{1/2}J$. Rango insuficiente produce infinito; las fronteras de corte se excluyen como \texttt{NaN}. Se exige regularidad local, no se derivan indicadores como si fueran funciones suaves. Una cota pequeña muy cerca del cero SQ puede ser optimista respecto a errores finitos y sensibilidad a un ajuste incorrecto.

\section{Identificabilidad: lo que sí es general}
Sea $E\in\mathbb R^{3\times2}$ una base tangente a la dirección. En coordenadas locales $(\log\beta,\boldsymbol\omega)$, el Jacobiano de medias es:
\begin{equation}
 L=\beta\,[\mathbf h,\ \operatorname{diag}(\mathbf h')HE],\qquad H_i=\mathbf n_i^{\mathsf T}.
\end{equation}
La transformación desde posición a esas coordenadas tiene matriz:
\begin{equation}
 A=\frac1d\begin{bmatrix}(2\mathbf u-\nabla_S\ln f_T)^{\mathsf T}\\-E^{\mathsf T}\end{bmatrix},
\end{equation}
invertible, pues su componente radial es $2/d$ y las otras dos son tangentes independientes. Así, en el dominio iluminado:
\begin{equation}
 \operatorname{rank}J=\operatorname{rank}[\mathbf h,\operatorname{diag}(\mathbf h')HE].
\label{eq:rank}
\end{equation}
Para el coseno-potencia esta condición se reduce al rango de normales visibles gracias a homogeneidad. Para una función general \textbf{no se debe reutilizar esa reducción}. Una respuesta afín/no homogénea puede aportar una relación de amplitud distinta. Se ha verificado un ejemplo SQ con normales contenidas en un plano ($\operatorname{rank}H=2$), Jacobiano de posición de rango tres y dos soluciones globales simétricas. Por ello identificabilidad local no implica unicidad global; los estimadores multisemilla sólo detectan algunas ambigüedades, no certifican su ausencia.

\section{Comprobación axial independiente}
Para un LED vertical, un PD en su eje y un cono uniforme de $K\geq3$ normales con inclinación $\theta$, todas activas, el patrón emisor es estacionario angularmente en el eje. Para $N_i=N$:
\begin{equation}
 I_r=\frac{NKC^2}{\sigma^2d^6}\operatorname{diag}\left(
 \frac{h'(\cos\theta)^2\sin^2\theta}{2},
 \frac{h'(\cos\theta)^2\sin^2\theta}{2},4h(\cos\theta)^2\right).
\end{equation}
De ahí:
\begin{equation}
 \mathrm{PEB}^2=\frac{\sigma^2d^6}{NKC^2}\left[
 \frac4{h'(\cos\theta)^2\sin^2\theta}+\frac1{4h(\cos\theta)^2}\right].
\end{equation}
Esta expresión prueba de forma independiente el código de FIM para las tres familias. En $\theta=0$ desaparece diversidad angular; no se convierte una cota singular en finita mediante pseudoinversa.

\section{Pose y parámetros de ajuste inciertos}
Con perturbación tangente $\delta_i$ de la normal, la fila de derivadas de pose es:
\begin{equation}
 D_i=\beta h'(s_i)\mathbf u^{\mathsf T}E_i.
\end{equation}
Si las medidas auxiliares de pose tienen covarianza $V_\theta$, el complemento de Schur da:
\begin{equation}
 I_e=J^{\mathsf T}(\Sigma_y+DV_\theta D^{\mathsf T})^{-1}J.
\end{equation}
Los experimentos K=5 comparan pose exacta con desviación estándar de $1^\circ$ por componente tangente. El error persiste sobre todas las muestras ópticas de su orientación. La covarianza equivalente no representa por sí sola una distribución marginal de jitter de actuación: aquí procede de eliminar poses de una verosimilitud conjunta con medidas auxiliares.

Los coeficientes publicados se tratan como conocidos en los resultados principales. Una incertidumbre en $\psi_{3\mathrm{dB}}$ requeriría otra columna auxiliar. Para SQ, dentro del soporte:
\begin{equation}
 \frac{\partial\mu_i}{\partial\psi_{3\mathrm{dB}}}
 =\frac{2\beta a\psi_i^2}{\psi_{3\mathrm{dB}}}.
\end{equation}
Se podría añadir información independiente de calibración y eliminar ese parámetro con otro complemento de Schur. El artículo no proporciona una covarianza de esos coeficientes que permita fijarla como dato medido. No se inventa una precisión de calibración a partir del número de decimales o de $R^2$.

\section{Valor para un journal y limitaciones}
\textbf{Es pertinente utilizar SQ como caso de estudio calibrado}, pero sustituir una función por otra no constituye por sí solo una contribución general fuerte. La formulación con $h(s)$ evita limitar la teoría a dos fotodiodos. Las contribuciones defendibles incluyen el criterio \eqref{eq:rank}, el diseño de reorientaciones bajo incertidumbre, la separación de soporte y pendiente, los estimadores para respuesta arbitraria y la evaluación de sesgo por desajuste.

Un buen ajuste de amplitud no asegura un buen ajuste de \emph{derivadas}. La FIM depende de $h'$; una caída idealizada muy abrupta puede producir información excesivamente favorable. Para afirmaciones de hardware se requiere caracterización de nuestra unidad, validación fuera de los puntos de ajuste, ruido adecuado y calibración del centro óptico, pose, ganancia y espectro. SQ predice soporte y respuesta de un modelo publicado; no acredita un error experimental real de Cambridge.

\section{Selección de modelos y lectura de los experimentos}
\subsection{Tres familias por defecto y cortes elegidos por el usuario}
Los scripts \path{sim01_paper_PD_responses} y \path{sim02_paper_K5_PEB} utilizan ahora, por defecto, los tres casos siguientes, en este orden:
\begin{enumerate}
\item \path{cosine_1}: $R(\psi)=\cos\psi$, con $m_R=1$.
\item \path{cosine}: $R(\psi)=\cos^{m_R}\psi$. El alias \path{cosine_mR} es equivalente.
\item \path{SQ}: la respuesta cuadrática publicada, con su cero natural y la apertura externa seleccionada.
\end{enumerate}
La lista \path{experiment.families} permite seleccionar y reordenar los modelos. SQapprox y Exp siguen disponibles si se solicitan explícitamente; ya no se añaden por defecto. Para comparar sólo coseno simple y SQ:
\begin{verbatim}
experiment.families = {'cosine_1','SQ'};
\end{verbatim}
La configuración predeterminada de tres modelos es:
\begin{verbatim}
experiment.families = {'cosine_1','cosine','SQ'};
experiment.cosine_fov_deg = 46;
experiment.outer_fov_deg = 90;
experiment.m_R = [];
\end{verbatim}
El primer corte es común a los dos cosenos. El segundo es la apertura externa de SQ y de las otras familias del artículo. Para SQ, el soporte efectivo es el mínimo entre su cero natural y ese corte externo. El valor vacío de $m_R$ utiliza el orden publicado del dispositivo: 1.9 para PDA100A2 y 0.98 para PDA36A2. Un valor explícito, por ejemplo 2.3, modifica sólo el coseno-potencia; \path{cosine_1} permanece en uno y SQ no utiliza ese exponente.

El corte receptor se configura independientemente de \path{transmitter.half_angle_power_deg}. Si el investigador llama $\Phi_{1/2}$ al ángulo que desea imponer como corte de recepción, debe introducir su valor en \path{cosine_fov_deg}; el programa no modifica por ello el patrón del LED. Tampoco deduce automáticamente un corte desde el nivel de media respuesta: el coseno simple alcanza 0.5 a $60^\circ$, y el de orden 1.9 a unos $46.03^\circ$. Aplicar 46 grados a ambos es una \emph{decisión de truncamiento común}, no afirmar que ambos tengan allí su media respuesta.

Las leyendas se generan a partir de los parámetros de cada curva. No deben reemplazarse con un comando \texttt{legend} independiente del orden de los datos. La tabla \path{models.csv} registra el identificador, el orden real, la FOV externa, el soporte efectivo y el documento de derivación aplicable. Los MAT históricos conservan los modelos con los que se calcularon: redibujarlos no los transforma en el nuevo conjunto de tres casos.

\subsection{Referencia explícita para el coseno simple}
El script K=5 expone \path{experiment.peb_reference_cosine_1}. Esta referencia apunta al documento original \path{PEB_derivation.tex}, dentro de la carpeta de cotas 3D. Es la derivación nominal para $m_R=1$ y normales conocidas. No se reemplaza ni modifica ese documento. La extensión a orientaciones inciertas utiliza, además, el modelo conjunto de RSS y medida de orientación descrito en \path{Reorientation sensitivity/Sensitivity_report.tex}. La columna \path{PEBDerivation} distingue la referencia original, la generalización coseno-potencia y la del presente documento.

\subsection{Qué representan los ejes}
\begin{description}
\item[Respuesta angular.] El eje horizontal es la incidencia $\psi$, en grados: el ángulo entre la normal del PD y la dirección hacia el LED. El vertical es $R(\psi)$, normalizado con $R(0)=1$, sin unidades. No representa vatios, error de posición ni porcentaje de cobertura.
\item[Derivada angular.] El eje horizontal sigue siendo $\psi$ en grados. El vertical es $|\mathrm dR/\mathrm d\psi|$ por \emph{radián}, pues las derivadas analíticas se calculan en radianes. Para expresarlo por grado se multiplicaría por $\pi/180$. La derivada regular no está definida en el corte; no se interpreta un salto como información infinita.
\item[Cobertura K=5.] El eje horizontal es el tilt común $\theta$ de las cinco normales, medido desde la vertical global. No es el ángulo de incidencia $\psi$, que también depende de la posición. El vertical es el porcentaje de \emph{todos} los puntos de la rejilla densa que tienen PEB finito y regular menor o igual al objetivo, por defecto 10 cm. Los infinitos y los NaN permanecen en el denominador como no cubiertos.
\item[RMS-PEB condicionado.] El eje horizontal es el mismo tilt $\theta$. El vertical, en cm, es la raíz del promedio de PEB al cuadrado sobre el subconjunto de puntos finitos. El subconjunto puede cambiar entre modelos e inclinaciones; esta curva debe leerse junto con la cobertura. No es RMSE medido de un estimador.
\item[Mapas espaciales.] Los ejes $x,y$ son coordenadas en metros. El valor $z$ del título identifica una sección horizontal del volumen; no se entrega como altura conocida al estimador. Los colores distinguen rango insuficiente, PEB finito por encima del objetivo, objetivo satisfecho y frontera no regular excluida. Se crea un mapa por modelo seleccionado.
\end{description}
Los ejes \emph{angulares} se muestran de 0 a 90 grados. Los ejes espaciales no se convierten a ese rango. La extensión visual a 90 grados no extrapola el barrido: con un límite mecánico configurado en 85 grados sólo se calculan inclinaciones hasta 85, y el tramo restante queda sin datos.

\subsection{Incertidumbre de orientación y actitud de referencia}
En los gráficos se utiliza $\sigma_{\mathrm{orient}}$, no el término ambiguo \emph{pose}. El parámetro \path{orientation_std_deg} es una desviación estándar por componente tangente local; no es una varianza ni un error máximo garantizado. No representa directamente una desviación estándar de azimut, especialmente cerca de inclinación cero.

El modo predeterminado \path{orientation_structure='independent'} considera conocida la referencia global y \emph{no añade un error común de actitud del soporte}. Sin embargo, las normales individuales del PD \emph{no son perfectamente conocidas} cuando $\sigma_{\mathrm{orient}}>0$: se dispone de medidas auxiliares de cada orientación con esa incertidumbre. Si por actitud se entiende la orientación absoluta del propio PD, sus dos componentes relevantes sí son inciertas; no son dos conceptos físicos independientes. El giro alrededor de la normal no interviene en una respuesta axisimétrica.

La cota es el complemento de Schur de la FIM conjunta RSS--orientación. El error permanece durante todas las muestras de un hold y no se divide por su número. Con $\sigma_{\mathrm{orient}}=0$ se recuperan las normales exactamente conocidas. Si se selecciona \path{common_rotation}, se introduce en cambio un error común de referencia de actitud, correlacionado entre adquisiciones, y los títulos usan $\sigma_{\mathrm{attitude}}$. Ninguno debe confundirse con el modelo de actuación aleatoria no observada ni con el PEB de normales perturbadas entregadas al estimador como exactas.

\section{Separación en el software}
\begin{itemize}
\item \path{System/Parameters/system_parameters.m}: conserva la familia coseno-potencia por defecto, con $m_R=1$ y la FOV configurada. El selector se llama \path{response_model}.
\item \path{Receiver models/rx_receiver_response.m}: despacho explícito de $h,h'$ y límites. Un MAT histórico sin selector sigue significando coseno-potencia.
\item \path{Receiver models/Bastiaens 2020/Parameters}: presets bibliográficos de SQ, SQapprox, Exp y coseno ajustado por dispositivo. Los presets no cambian área, ruido ni electrónica implícitamente.
\item \path{System/rx_channel.m} y \path{rx_peb.m}: física y FIM comunes con derivadas acordes al modelo seleccionado.
\item \path{Estimators}: los estimadores históricos de raíces coseno rechazan explícitamente un perfil SQ/Exp. No usan un $m_R=1$ ficticio.
\item \path{Receiver models/Bastiaens 2020/Estimators}: algoritmos propios para respuesta calibrada general, explicados en \path{PD_estimators.tex}.
\end{itemize}
Los campos $m_R$ se retiran al crear un preset no coseno para evitar parámetros contradictorios. La descripción de modelo impresa y guardada contiene familia, dispositivo, coeficientes activos, FOV externa, soporte efectivo y área realmente utilizada. El informe \path{PD_K5_results.tex} presenta los resultados y comandos de los experimentos independientes.

\bibliographystyle{IEEEtran}
\bibliography{bastiaens2020_references}
\end{document}
